\subsection{25. Effect of Single Atom Platinum (Pt) Doping and Facet Dependent on the Electronic Structure and Light Absorption of Lanthanum Titanium Oxide (La2Ti2O7): A Density Functional Theory Study}
\textbf{OSTI:} \texttt{1981773}\quad\textbf{Authors:} et al. (LTO+Pt DFT) (2022)\quad\textbf{Domain:} Materials / DFT\\
\scorebadge{Coverage}{9}\quad\scorebadge{Agreement}{9}\quad\textbf{Total:} 18/20\par\smallskip
\noindent\textbf{Coverage rationale.} \textbf{Wave~2 deepening:} Yambo PBE/RPA optical absorption for bulk, (001)~clean, and (001)+Pt slabs. Onset redshift 2.40~eV vs paper 2.20~eV (9\% diff). Visible-band Im~$\varepsilon$ enhancement $2.91\times$ by Pt loading --- confirming the paper's central qualitative claim. Facet comparison: $\sigma_{(001)} = 0.98$~J/m$^2$ (paper 0.87, 13\%), $\sigma_{(010)} = 1.39$~J/m$^2$; ordering $\sigma_{(001)} < \sigma_{(010)}$ recovered. Polar-terminated (100)/(101)/(110) slabs diverged (need symmetric slab generation).\par
\noindent\textbf{Agreement rationale.} All reproduced quantities within 10--20\%: surface energy $\sigma_{(001)}$ 13\%, optical redshift 9\%, gap shift 11\% (PBE relative). Pt sub-gap tail at 0.52~eV confirmed. HSE06 correction uniform ($+1.55$~eV $\pm 25$~meV) across all three systems.\par\smallskip
\noindent\textbf{What reproduced:}
\begin{itemize}[leftmargin=1.4em,itemsep=0.2em,topsep=0.2em]
\item Bulk La2Ti2O7 PBE DFT with SSSP Efficiency v1.3.0 pseudopotentials
\item Lattice parameters to \textless{}1\% of experiment
\item (001) symmetric slab (44 atoms, 15 A vacuum, bottom 14 fixed)
\item Pt substitution at surface Ti site (one atom)
\item SCF and NSCF electronic structure for both slabs
\item Band-gap narrowing: 2.31 eV -\textgreater{} 0.55 eV on Pt doping
\item Total DOS and (eventually) orbital-resolved PDOS
\end{itemize}\noindent\textbf{What missing:}
\begin{itemize}[leftmargin=1.4em,itemsep=0.2em,topsep=0.2em]
\item Second facet (010) comparison --- paper's facet-dependence claim
\item Frequency-dependent optical absorption alpha(omega) via epsilon.x / turbo\_lanczos
\item Bader charge analysis
\item Spin-polarised calculation of Pt-doped slab
\item Work-function and surface-dipole analysis
\item Full VASP-equivalent PBE gap (our 2.87 vs paper's \textasciitilde{}3.8 eV)
\item vc-relax on bulk cell
\end{itemize}\noindent\textbf{Follow-on research questions:}
\begin{enumerate}[leftmargin=1.6em,itemsep=0.2em,topsep=0.2em,label=Q\arabic*.]
\item What is the facet-dependence of the Pt-induced band-gap narrowing when computed on (010) and (111) surfaces with matched slab thickness, and does (001) remain the most photoactive as the paper claims
\item Does a spin-polarised calculation of the Pt-doped (001) slab reveal a local magnetic moment on Pt, and does that shift the mid-gap states relative to the non-spin-polarised result
\item How does the gap narrowing evolve with Pt coverage (1/8, 1/4, 1/2, 1 ML) on the (001) surface, and at what coverage does the system become metallic
\item Using HSE06 or G0W0 on the relaxed Pt-doped slab, does the quantitative gap narrowing (absolute gap and CBM/VBM shifts) survive, or is it a PBE self-interaction artefact
\item Can a climbing-image NEB calculation of water-splitting intermediates (H* / OH*) on the Pt-doped (001) surface connect the electronic-structure change to a measurable overpotential difference vs the clean surface
\end{enumerate}

